% Generated draft: requires mathematical and source review.
\documentclass[11pt]{article}
\usepackage[margin=25mm]{geometry}
\usepackage{fontspec}
\setmainfont{Latin Modern Roman}
\usepackage{amsmath,amssymb,mathtools,unicode-math}
\setmathfont{Latin Modern Math}
\usepackage{enumitem,needspace,xurl,hyperref}
\setlength{\parindent}{0pt}
\setlength{\parskip}{5pt}
\setlength{\emergencystretch}{3em}
\newcommand{\N}{\mathbb{N}}
\newcommand{\Z}{\mathbb{Z}}
\newcommand{\Q}{\mathbb{Q}}
\newcommand{\R}{\mathbb{R}}
\newcommand{\C}{\mathbb{C}}
\newcommand{\F}{\mathbb{F}}
\newcommand{\A}{\mathbb{A}}
\newcommand{\PP}{\mathbb{P}}
\newcommand{\E}{\mathbb{E}}
\newcommand{\eps}{\varepsilon}
\newcommand{\dd}{\,\mathrm{d}}
\begin{document}
\section*{UnsolvedMath ID: 30004448}
\textbf{Source problem code:} catalogue entry
\subsection{Definitions and mathematical statement}
A Polish group is a topological group whose topology is separable and induced by a complete metric. Here separability means the existence of a countable dense subset. An abstract group homomorphism \(f:G\to H\) satisfies
\[
 f(xy)=f(x)f(y)\qquad(x,y\in G),
\]
without an assumed continuity or measurability condition. Write \(e_G\) and \(e_H\) for the identity elements. Such a homomorphism is discontinuous precisely when
\[
\begin{gathered}
\exists V\subseteq H\;\bigl[V\text{ is open},\ e_H\in V,\\
\forall U\subseteq G\;\bigl((U\text{ is open and }e_G\in U)
\Longrightarrow \exists x\in U\; f(x)\notin V\bigr)\bigr].
\end{gathered}
\]
This uses the equivalence between continuity everywhere and continuity at the identity for group homomorphisms.

The supplied question is: \emph{How much choice is needed to construct a discontinuous homomorphism between Polish groups?}

The existential assertion naturally associated with this question is
\[
\exists G\,\exists H\,\exists f\;
\bigl(G,H\text{ are Polish groups and }f:G\to H
\text{ is a discontinuous group homomorphism}\bigr).
\]
The problem asks for the choice principles necessary or sufficient for such a construction. The source does not specify a background set theory, a class or ordering of candidate choice principles, or whether the groups are to be fixed beforehand. Consequently, the displayed assertion records an existential interpretation; it does not resolve whether the intended target is an equivalence over a particular base theory, a relative consistency comparison, or a construction for specified groups. These qualifications must be settled before a unique formal statement of minimal choice strength can be given.
\subsection{Short English statement}
Determine the strength of choice needed to obtain a discontinuous abstract homomorphism between Polish groups. The source leaves the background theory, the meaning of minimal choice strength, and the intended quantification over groups unspecified. Its supplied status is partially solved; this draft does not independently verify that assessment.
\subsection{Sources}
\begin{itemize}
\item Supplied local UnsolvedMath registry record, permanent ID \(30004448\), title \emph{Choice Principles for Discontinuous Polish Group Homomorphisms}, source problem code \texttt{catalogue entry}. The record gives proposed year \(2020\), marks the statement as exact, and reports that its wording was retained byte-for-byte. No proposer is identified.
\item Supplied report citation: \emph{Set Theory (2021)}, Oberwolfach Reports, pp.\ 797--855, DOI \texttt{10.4171/owr/2020/14}. Supplied source URL and literature reference labeled \emph{Oberwolfach Report 14/2020}: \url{https://doi.org/10.4171/owr/2020/14}. The supplied background describes these reports as documenting mathematical workshops and their open-problem discussions.
\item Supplied literature assessment, checked on \textbf{2026-08-22}: status \emph{partially solved}. The curation attributes strong continuity results for homomorphisms of Polish groups and lower bounds on the choice needed for discontinuous homomorphisms to Rosendal, while stating that the exact minimal choice principle is not determined by the reported results. This is supplied dated curation, not an independently verified literature assessment. The record explicitly requests independent verification of time-sensitive status and cited claims.
\end{itemize}
\end{document}
